\vfill\pagebreak
\section{Where values are kept}
\label{sec:mutability:stack_heap}

A running program keeps its values in the memory of the machine, a long row of
bytes, each with a number called its \textit{address}. The program uses two
regions of it for the values of its variables.

\begin{itemize}
\item \textit{The stack.} Each call of a function takes a zone of the stack for
  its parameters and its variables, called its \textit{frame}. The frame is
  taken when the call starts and given back when it returns, so the variables
  of a function exist only during its call. A call made from inside the
  function gets a frame of its own, next to that of its caller: this is the
  call stack of Section~\ref{sec:functions:infinite_recursion}, which
  overflows when too many calls wait at the same time.
\item \textit{The heap.} The program can also reserve a block of memory
  anywhere in the heap, of any size, when it runs. Chapter~\ref{chap:collections}
  called that an allocation. A block does not belong to a call: it stays as long
  as the program needs it, whichever functions return in the meantime.
\end{itemize}

The next listing declares an array and a slice, and
Figure~\ref{fig:mutability:array_slice} shows where their elements are.

\lstinputlisting[style=coloredverbatim, caption={An array and a slice, \textit{examples/mutability/where.yr}}, label=lst:mutability:where]{examples/mutability/where.yr}
\vspace{-10pt}%
\begin{lstlisting}[style=bashVerb]
[31, 28, 31] [2, 3, 5, 7]
\end{lstlisting}

\begin{figure}[h]
  \centering
  \begin{adjustbox}{max width=\linewidth}
    \begin{tikzpicture}
      \foreach \value [count=\i from 0] in {31, 28, 31} {
        \node[memcell] (d\i) at (\i * 1.0, 0) {\value};
      }
      \node[cflab, left=3mm of d0] (dn) {\token{days}};
      \node[memcell] (len) at (0, -1.2) {4};
      \node[memaddr, right=0mm of len] (ptr) {};
      \node[cflab, left=3mm of len] (pn) {\token{primes}};
      \node[memnote, below=0.5mm of len] {length};
      \node[memnote, below=0.5mm of ptr] {address};
      \node[memframe, fit=(dn)(d2)(pn)(ptr), inner ysep=4mm] (main) {};
      \node[memtitle] at ([xshift=2mm]main.north west) {main};

      \foreach \value [count=\i from 0] in {2, 3, 5, 7} {
        \node[memcell] (h\i) at (5.2 + \i * 1.0, -1.2) {\value};
      }
      \draw[mempoint] (ptr.center) -- (h0.west);

      \node[memregion] (stack) at ($(main.north) + (0, 5mm)$) {stack};
      \node[memregion] at ({$(h0)!0.5!(h3)$} |- stack.base) {heap};
      \draw[memsep] (4.3, 0.9) -- (4.3, -2.4);
    \end{tikzpicture}
  \end{adjustbox}
  \caption{\label{fig:mutability:array_slice} The array \token{days} and the
    slice \token{primes} of Listing~\ref{lst:mutability:where}. The elements of
    the array are in the variable, in the frame of \token{main} on the stack.
    The slice holds a length and an address, and its elements are in a block
    of the heap.}
\end{figure}


\subsection{Values stored in the variable}

The three elements of \token{days} are in the variable itself, one after the
other, in the frame of \token{main}. The same goes for every type met so far
except slices and maps: an integer, a float, a character or a boolean take a few
bytes of the frame, and a tuple or an array take the bytes of all their fields
or elements. The compiler knows that size, since the length of an array is part
of its type, and reserves the frame accordingly.

\subsection{Values stored elsewhere}

A slice cannot be stored that way: its length is known only when the program
runs, and it can grow. The variable \token{primes} holds two numbers of a fixed
size instead: the length of the slice, and the address of its first element.
The elements are in a block of the heap, one after the other, as those of an
array are. Reading \token{primes[2]} follows the address, and counts two
elements from there. (A slice keeps a third number, used when it grows, which
the figures of this chapter leave out.)

A value whose elements lie outside of it is said to \textit{borrow} them: the
elements are not part of the value, the value only says where they are. Slices,
strings, which are slices of characters, and maps borrow their elements. A map
is laid out differently from a slice, but the principle is the same: its
variable holds the address of a structure of the heap, where the entries are.

The empty slice \token{[]} has a length of 0 and no element to designate, so it
needs no block, and no \token{copy}. A map refers to its structure even when it
has no entry, and \token{copy []} is what allocates that structure, which is
why an empty map cannot do without it (Section~\ref{sec:collections:maps}).

This answers a question of Chapter~\ref{chap:collections}. \token{[2, 3, 5, 7]}
is an array: its four elements are the value itself, and would be kept in the
frame. \token{copy} reserves a block of the heap, copies the four elements into
it, and gives a slice that borrows them. The block outlives the call that made
it, which is why a function can return a slice that it created.

A block of the heap is never given back by the program itself. The
\textit{garbage collector}, a part of every Ymir program, finds the blocks that
no value refers to any more, and reuses their memory. A program never has to
free a block, and never reads one that was freed.

\notebox{A string literal such as \token{"hello"} is a slice too, but its
  characters are not in the heap: the compiler writes them into the program
  itself, in a region that cannot change. \token{copy "hello"} copies them into
  a block of the heap, where they can change.}
